% Generated from locale/tr/content/normal-modal-logic/completeness/introduction.tex; do not hand-edit.


\olfileid[tr]{nml}{com}{int}

\olsection{Giriş}

$\Sigma$ bir modal sistemse, sağlamlık teoremi şunu kurar: $\Sigma\Proves !A$
ise $!A$, $\Sigma$ içindeki bütün !!{formula}s{}in bütün örneklerinin geçerli
olduğu her $\mClass{C}$ modeller sınıfında geçerlidir. Özellikle
$\Log{K}\Proves !A$ ise $!A$ bütün modellerde; $\Log{KT}\Proves !A$ ise
bütün yansımalı modellerde; $\Log{KD}\Proves !A$ ise bütün seri modellerde
doğrudur ve böyle devam eder.

Tamlık sağlamlığın tersidir: $\Log{K}$'nin tam olması, !!a{formula} $!A$
geçerliyse $\Log{K}\Proves !A$ olması demektir. Tamlığı ispatlamak,
sağlamlığı ispatlamaktan çok daha zordur. Önce karşıt-ters ifadeyi ele almak
yararlıdır: $\Log{K}$ ancak ve ancak $\Log{K}\Proves/ !A$ olduğunda bir
karşı-model, yani $\mSat/{M}{!A}$ olacak bir $\mModel{M}$ modeli varsa
tamdır. Denk olarak ($!A$'yı değilleyerek), $\Log{K}\Proves/\lnot !A$
olduğunda $!A$'nın bir modeli bulunduğunu ispatlayabiliriz. Böyle bir modeli
kurarken $!A$ içindeki bilgiyi kullanabiliriz. Belirli !!{formula}s için
model ararken de çoğu zaman aynısını yaparız: örneğin $p\lif\Box q$ için bir
karşı-model bulmak istiyorsak, modelde $p$'nin doğru ve $\Box q$'nun yanlış
olduğu bir dünya bulunması gerektiğini biliriz. $\Box q$'nun bir dünyada
yanlış olması, o dünyadan erişilebilen ve $q$'nun yanlış olduğu bir dünya
bulunması demektir. Bilmemiz gerekenlerin tümü budur: hangi dünyaların hangi
!!{propositional variable}s{}ni doğru yaptığı ve hangi dünyalara hangi
dünyalardan erişilebildiği.

Ne var ki tamlığı ispatlarken kendisi için model kurduğumuz belirli bir
!!{formula} $!A$ yoktur. $\Proves/[\Sigma]\lnot !A$ olan her $!A$ için bir
model bulunduğunu kurmak isteriz. Bu asgari bir koşuldur; çünkü
\emph{eğer} $\Proves[\Sigma]\lnot !A$ ise sağlamlık gereğince $!A$'nın
($\Sigma$'nın doğru olduğu) bir modeli yoktur. Şimdi
$\Proves/[\Sigma]\lnot !A$ ancak ve ancak $!A$ $\Sigma$-tutarlıysa
geçerlidir. ($\Proves/[\Sigma]\lnot !A$ ile
$!A\Proves/[\Sigma]\lfalse$ ifadelerinin denk olduğunu hatırlayın.)
Dolayısıyla görevimiz her $\Sigma$-tutarlı !!{formula} için bir model
kurmaktır.

Kullanacağımız yöntem, $!A$'yı içeren ve aynı zamanda $!A$'yı doğru yapan
dünyanın nasıl olması gerektiğini bildiren başka formülleri de içeren
$\Sigma$-tutarlı bir !!{formula}s kümesi bulmaktır. Böyle kümeler
\emph{tam} $\Sigma$-tutarlı kümelerdir. $!A$'yı doğru yapmak için tek
dünyalı bir model kurmak yeterli değildir; model birden fazla dünya ve bir
erişilebilirlik bağıntısı içermelidir. $!A$'yı içeren tam $\Sigma$-tutarlı
küme, $\Box !B$ ve $\Diamond !C$ biçiminde başka !!{formula}s da içerir.
Bütün erişilebilir dünyalarda $!B$; en az birinde $!C$ doğru olmalıdır.
Bunu sağlamak için \emph{bütün} olası tam $\Sigma$-tutarlı kümeleri dünyalar
kümesinin temeli olarak alacağız. Zor kısım, modelimizde bir tam
$\Sigma$-tutarlı kümenin başka birinden ne zaman erişilebilir sayılacağını
belirlemektir.

Bu biçimde tanımlanan modelde $!A$'nın bir dünyada---kendisi de tam
$\Sigma$-tutarlı bir kümedir---ancak ve ancak $!A$ o kümenin !!a{element}{}ı
ise doğru olduğunu göstereceğiz. $!A$ $\Sigma$-tutarlıysa en az bir tam
$\Sigma$-tutarlı kümenin !!a{element}{}ı olacaktır (bunu ispatlayacağız) ve
dolayısıyla $!A$'nın doğru olduğu bir dünya bulunacaktır. Böylece her
$\Sigma$-tutarlı !!{formula} $!A$'nın bir dünyada doğru olduğu tek bir
model elde edeceğiz. Bu tek model, $\Sigma$ için \emph{kanonik} modeldir.

